\documentclass[12pt]{article}
\usepackage[utf8]{inputenc}
\usepackage[T1]{fontenc}
\usepackage[margin=1in]{geometry}
\usepackage{setspace}
\usepackage{xcolor}
\usepackage[colorlinks=true, urlcolor=blue!60!black, linkcolor=blue!60!black]{hyperref}
\setlength{\parindent}{2em}
\setlength{\parskip}{1.25em}
\renewcommand{\baselinestretch}{1.0}
\begin{document}
\begin{center}
{\LARGE \textit{United States v. Texas}, 599 U.S. 670 (2023)\par}
\vspace{0.8em}
{\large Case Brief\par}
\end{center}
\vspace{1em}

\section{Memory Jogger}

States lack Article III standing to sue the federal Executive Branch to compel more arrests or prosecutions of third-party noncitizens, because indirect fiscal costs from executive non-enforcement are not a traditionally judicially cognizable injury (\href{https://www.courtlistener.com/opinion/10049663/united-states-v-texas/}{CourtListener}; \href{https://www.law.cornell.edu/supremecourt/text/22-58}{Cornell LII}).

\section{Detailed Case Facts}

Because there are more than 11 million removable noncitizens in the United States and the Department of Homeland Security (DHS) lacks resources to arrest and remove all of them, presidential administrations have historically set civil immigration enforcement priorities (\href{https://www.law.cornell.edu/supremecourt/text/22-58}{Cornell LII}; \href{https://www.lawfaremedia.org/article/united-states-v.-texas-explained}{Lawfare}). On September 30, 2021, following interim memoranda issued earlier that year, DHS Secretary Alejandro Mayorkas promulgated the ``Guidelines for the Enforcement of Civil Immigration Law'' (\href{https://harvardlawreview.org/print/vol-137/united-states-v-texas/}{Harvard Law Review}; \href{https://www.law.cornell.edu/supremecourt/text/22-58}{Cornell LII}). The Guidelines prioritized the apprehension and removal of noncitizens who pose a threat to national security (such as suspected terrorists or spies), public safety (dangerous criminals), or border security (recent unlawful entrants), and instructed DHS personnel to evaluate individualized aggravating and mitigating factors rather than relying solely on the fact of a criminal conviction (\href{https://www.law.cornell.edu/supremecourt/text/22-58}{Cornell LII}; \href{https://www.lawfaremedia.org/article/united-states-v.-texas-explained}{Lawfare}).

Texas and Louisiana sued DHS and federal officials under the Administrative Procedure Act (APA), 5 U.S.C. \S 706(2), alleging that the Guidelines violated two provisions of the Immigration and Nationality Act (INA) enacted in the Illegal Immigration Reform and Immigrant Responsibility Act of 1996 (IIRIRA): (1) 8 U.S.C. \S 1226(c), which states that the Attorney General ``shall take into custody'' certain criminal and terrorist noncitizens when released from state criminal custody; and (2) 8 U.S.C. \S 1231(a)(2), which states that the government ``shall detain'' noncitizens during the 90-day removal period after a final removal order (\href{https://www.law.cornell.edu/supremecourt/text/22-58}{Cornell LII}; \href{https://lawreview.law.lsu.edu/archives/3081}{Louisiana Law Review}).

The States contended that DHS's failure to arrest and take custody of all covered noncitizens forced Texas and Louisiana to bear four categories of ongoing financial costs: continued state criminal supervision (probation, parole, or incarceration) after Immigration and Customs Enforcement (ICE) rescinded detainers, costs of criminal recidivism, public education expenses for un-detained juvenile offenders, and state-funded healthcare expenses (\href{https://www.law.cornell.edu/supremecourt/text/22-58}{Cornell LII}; \href{https://lawreview.law.lsu.edu/archives/3081}{Louisiana Law Review}).

\section{Procedural History}

\begin{itemize}
\item \textbf{U.S. District Court for the Southern District of Texas:} Following a two-day bench trial, Judge Drew Tipton found that the Guidelines caused ICE to rescind detainers and imposed concrete monetary costs on Texas, conferring Article III standing. On June 10, 2022, in \textit{Texas v. United States}, 606 F. Supp. 3d 437 (S.D. Tex. 2022), the court held on the merits that the Guidelines were contrary to law, arbitrary and capricious, and issued without notice-and-comment rulemaking under the APA, and vacated the Guidelines nationwide under 5 U.S.C. \S 706(2) (declining to issue an injunction or declaratory judgment) (\href{https://www.law.cornell.edu/supremecourt/text/22-58}{Cornell LII}; \href{https://ballotpedia.org/United_States_v._Texas_(2023)}{Ballotpedia}).
\item \textbf{U.S. Court of Appeals for the Fifth Circuit:} The Fifth Circuit denied the federal government's motion for a stay pending appeal in \textit{Texas v. United States}, 40 F.4th 205 (5th Cir. 2022), concluding that the government had not shown clear error in the district court's standing findings and was unlikely to succeed on the merits (\href{https://www.law.cornell.edu/supremecourt/text/22-58}{Cornell LII}; \href{https://harvardlawreview.org/print/vol-137/united-states-v-texas/}{Harvard Law Review}).
\item \textbf{U.S. Supreme Court:} On July 8, 2022, the United States applied for a stay and petitioned for certiorari before judgment. On July 21, 2022, the Supreme Court denied the stay (5--4) but granted certiorari before judgment on three questions: Article III standing, compliance with 8 U.S.C. \S\S 1226(c) and 1231(a) and the APA, and whether 8 U.S.C. \S 1252(f)(1) bars APA vacatur under 5 U.S.C. \S 706(2) (\href{https://www.scotusblog.com/cases/united-states-v-texas-5/}{SCOTUSblog}; \href{https://ballotpedia.org/United_States_v._Texas_(2023)}{Ballotpedia}). The Court heard oral argument on November 29, 2022, and reversed on June 23, 2023, resolving the case solely on Article III standing (\href{https://www.law.cornell.edu/supremecourt/text/22-58}{Cornell LII}).
\end{itemize}

\section{Judicial Votes}

\textbf{8--1 decision} reversing the district court (\href{https://www.law.cornell.edu/supremecourt/text/22-58}{Cornell LII}; \href{https://www.scotusblog.com/cases/united-states-v-texas-5/}{SCOTUSblog}):

\begin{itemize}
\item \textbf{Majority Opinion (Justice Kavanaugh):} Joined by Chief Justice Roberts and Justices Sotomayor, Kagan, and Jackson (5 Justices).
\item \textbf{Concurring in the Judgment (Justice Gorsuch):} Joined by Justices Thomas and Barrett (3 Justices).
\item \textbf{Concurring in the Judgment (Justice Barrett):} Joined by Justice Gorsuch (2 Justices).
\item \textbf{Dissenting Opinion (Justice Alito):} Sole dissenter (1 Justice).
\end{itemize}

\section{Holding}

\textbf{Judgment:} Reversed (8--1). Because Texas and Louisiana lacked Article III standing, the U.S. District Court for the Southern District of Texas lacked subject-matter jurisdiction (\href{https://www.law.cornell.edu/supremecourt/text/22-58}{Cornell LII}).

\textbf{Rule Announced:} Under Article III of the Constitution, a plaintiff---including a State---lacks a legally and judicially cognizable interest in challenging the Executive Branch's arrest or prosecution policies when the plaintiff is neither prosecuted nor threatened with prosecution, even if executive under-enforcement indirectly imposes monetary costs on the plaintiff (\href{https://www.law.cornell.edu/supremecourt/text/22-58}{Cornell LII}; \href{https://supreme.justia.com/cases/federal/us/599/22-58/}{Justia}).

\textbf{Holding vs. Dicta:}
\begin{itemize}
\item \textbf{Holding:} Texas and Louisiana lack Article III standing to challenge DHS's 2021 civil immigration enforcement Guidelines because federal courts have not traditionally entertained suits ordering the Executive Branch to make more arrests or bring more prosecutions (\href{https://www.law.cornell.edu/supremecourt/text/22-58}{Cornell LII}).
\item \textbf{Dicta / Caveats Left Open:} The Court did \textit{not} decide whether DHS's Guidelines comply with 8 U.S.C. \S\S 1226(c) and 1231(a)(2), nor did it decide whether 8 U.S.C. \S 1252(f)(1) or 5 U.S.C. \S 706(2) permits universal APA vacatur. In Part II-C, the majority listed five situations in which the standing analysis ``might'' or ``arguably could'' differ---selective-prosecution Equal Protection claims, express congressional elevation of de facto injuries with explicit judicial remedies, complete executive ``abdication'' of statutory enforcement duties, policies combining non-enforcement with the conferral of legal benefits or legal status (such as work authorization), and policies governing continued detention of already-arrested noncitizens---without definitively resolving those hypotheticals (\href{https://www.law.cornell.edu/supremecourt/text/22-58}{Cornell LII}).
\end{itemize}

\section{Analysis of Opinions}

\subsection{Majority Opinion (Justice Kavanaugh)}

Writing for a five-Justice majority, Justice Kavanaugh grounded Article III standing in the separation of powers and historical practice (\href{https://www.law.cornell.edu/supremecourt/text/22-58}{Cornell LII}). Although monetary costs are an injury in fact, \textit{Raines v. Byrd}, 521 U.S. 811 (1997), and \textit{TransUnion LLC v. Ramirez}, 594 U.S. 413 (2021), require that an asserted injury also be ``legally and judicially cognizable,'' meaning the dispute is traditionally capable of judicial resolution. The States cited no precedent, history, or tradition of federal courts ordering the Executive Branch to alter its arrest policies to make more arrests. To the contrary, \textit{Linda R. S. v. Richard D.}, 410 U.S. 614, 619 (1973), established that ``a citizen lacks standing to contest the policies of the prosecuting authority when he himself is neither prosecuted nor threatened with prosecution'' (\href{https://www.law.cornell.edu/supremecourt/text/22-58}{Cornell LII}).

Kavanaugh identified three structural reasons supporting this rule. First, declining to arrest or prosecute a third party does not exercise coercive power over the plaintiff's liberty or property. Second, suits compelling arrests intrude on the Executive's Article II authority to prioritize law enforcement---an authority that extends to immigration, where enforcement implicates foreign-policy objectives (\textit{Reno v. American-Arab Anti-Discrimination Committee}, 525 U.S. 471 (1999); \textit{Arizona v. United States}, 567 U.S. 387 (2012)). Third, because the Executive always faces resource constraints (as five consecutive administrations recognized under \S\S 1226(c) and 1231(a)(2)), federal courts lack meaningful standards to second-guess enforcement trade-offs (\textit{Heckler v. Chaney}, 470 U.S. 821 (1985)) (\href{https://www.law.cornell.edu/supremecourt/text/22-58}{Cornell LII}).

Kavanaugh emphasized the narrowness of the holding by enumerating five potential exceptions (Part II-C), distinguishing \textit{Massachusetts v. EPA}, 549 U.S. 497 (2007), in footnote 6 as a rulemaking-petition case rather than an enforcement-discretion challenge, and stressing that ``the question of reviewability is different from the question of legality'': Congress and voters retain political checks such as appropriations, oversight, Senate confirmations, and elections (\href{https://www.law.cornell.edu/supremecourt/text/22-58}{Cornell LII}).

\subsection{Concurring in the Judgment (Justice Gorsuch, Joined by Justices Thomas and Barrett)}

Justice Gorsuch agreed that the States lacked standing, but diagnosed the defect as a lack of \textit{redressability} rather than a lack of a judicially cognizable injury (\href{https://www.law.cornell.edu/supremecourt/text/22-58}{Cornell LII}). Gorsuch questioned the majority's dismissal of the States' proven monetary costs, noted that ``special solicitude'' from \textit{Massachusetts v. EPA} had no basis in Article III and should be left ``on the shelf,'' and doubted the majority's attempt to cabin Article II enforcement discretion solely to arrests and prosecutions.

Instead, Gorsuch reasoned that 8 U.S.C. \S 1252(f)(1), as construed in \textit{Garland v. Aleman Gonzalez}, 596 U.S. 543 (2022), barred lower federal courts from enjoining federal immigration enforcement. Moreover, vacating the Guidelines under 5 U.S.C. \S 706(2) could not redress the States' injuries because DHS officials would retain the exact same underlying prosecutorial discretion without the Guidelines, and redressability turns on a court's binding judgment rather than the persuasive force of its opinion (\textit{Haaland v. Brackeen}, 599 U.S. 255 (2023)) (\href{https://www.law.cornell.edu/supremecourt/text/22-58}{Cornell LII}). In Part III, Gorsuch launched an extensive textual and historical critique of universal APA vacatur, arguing that ``set aside'' in \S 706(2) likely means ``disregard'' an unlawful rule as to the parties before the court---not erase it nationwide---and warning that universal vacatur replicates the separation-of-powers abuses of universal injunctions (\href{https://www.law.cornell.edu/supremecourt/text/22-58}{Cornell LII}; \href{https://www.lawfaremedia.org/article/united-states-v.-texas-explained}{Lawfare}).

\subsection{Concurring in the Judgment (Justice Barrett, Joined by Justice Gorsuch)}

Justice Barrett joined Justice Gorsuch's redressability rationale and wrote separately to critique the majority's doctrinal foundation (\href{https://www.law.cornell.edu/supremecourt/text/22-58}{Cornell LII}). First, Barrett argued that the majority overread \textit{Linda R. S. v. Richard D.}, which actually turned on speculative redressability (jailing a father would not necessarily produce child-support payments), and was inapposite because Texas and Louisiana did not demand the prosecution of any specific individual. Second, Barrett objected that the majority conflated merits and statutory-review doctrines---specifically \textit{Town of Castle Rock v. Gonzales}, 545 U.S. 748 (2005) (a Fourteenth Amendment due process case), \textit{Heckler v. Chaney} (an APA statutory reviewability case), and the Article II Take Care Clause---with Article III's constitutional definition of cognizable injury. Elevating statutory doctrines like \textit{Heckler} into Article III standing rules, she warned, creates a novel constitutional barrier beyond Congress's reach (\href{https://www.law.cornell.edu/supremecourt/text/22-58}{Cornell LII}).

\subsection{Dissenting Opinion (Justice Alito)}

Dissenting alone, Justice Alito argued that Texas straightforwardly satisfied the three-part \textit{Lujan v. Defenders of Wildlife} test: the district court made unchallenged factual findings after a bench trial that the Guidelines caused ICE to rescind detainers and forced Texas to incur concrete supervision, recidivism, education, and healthcare costs, which vacatur would redress (\href{https://www.law.cornell.edu/supremecourt/text/22-58}{Cornell LII}). Alito emphasized that Congress enacted 8 U.S.C. \S\S 1226(c) and 1231(a)(2) precisely to strip the Executive of discretion over criminal noncitizens, and prior precedents (\textit{Nielsen v. Preap}, 586 U.S. 392 (2019); \textit{Johnson v. Guzman Chavez}, 594 U.S. 523 (2021)) confirmed that ``shall'' means ``must.''

Alito accused the majority of brushing aside \textit{Massachusetts v. EPA}---noting that Texas surrendered its sovereign border-control powers upon entering the Union just as Massachusetts surrendered its sovereign prerogatives---and ignoring standing implicitly accepted in \textit{Biden v. Texas}, 597 U.S. 785 (2022), and \textit{Little Sisters of the Poor v. Pennsylvania}, 591 U.S. 657 (2020). He dismissed the majority's five caveats as vague, unprincipled, or incoherent. Most fundamentally, Alito argued that insulating executive disobedience of mandatory statutes from judicial review revives the pre-1688 royal suspending and dispensing powers repudiated by the English Bill of Rights of 1689, early state constitutions, the Take Care Clause, and \textit{Kendall v. United States ex rel. Stokes}, 37 U.S. (12 Pet.) 524 (1838), forcing Congress into ``inter-branch warfare'' to compel compliance with the law (\href{https://www.law.cornell.edu/supremecourt/text/22-58}{Cornell LII}).

\section{Examples: Future Applications}

\begin{enumerate}
\item \textbf{Same-Side (No Standing):} A state sues the U.S. Department of Justice challenging a federal memorandum that deprioritizes federal prosecutions of simple marijuana possession or low-level firearms offenses, claiming the policy increases state policing and incarceration costs. Under \textit{United States v. Texas}, the state lacks Article III standing because it is neither prosecuted nor threatened with prosecution and seeks to compel the Executive Branch to make more arrests and bring more prosecutions (\href{https://www.law.cornell.edu/supremecourt/text/22-58}{Cornell LII}).
\item \textbf{Same-Side (No Standing):} A coalition of municipal police unions and border counties sues the Bureau of Alcohol, Tobacco, Firearms and Explosives (ATF) alleging that the agency's guidance prioritizing violent-crime investigations over technical paperwork violations violates a federal statute stating the agency ``shall'' investigate all reported violations, thereby increasing local law-enforcement expenditures. As in \textit{Texas}, indirect fiscal burdens from executive enforcement prioritization are not a judicially cognizable injury under Article III (\href{https://www.law.cornell.edu/supremecourt/text/22-58}{Cornell LII}).
\item \textbf{Opposite-Side (Standing Exists):} A state sues DHS to challenge a federal program that not only defers deportation for a class of noncitizens, but also affirmatively confers lawful presence, work authorization, and eligibility for state-subsidized driver's licenses or public benefits (similar to the DAPA program). Under the majority's fourth caveat in Part II-C, a policy combining non-enforcement with affirmative legal benefits or legal status implicates more than pure arrest discretion and can support state standing based on direct fiscal outlays (\href{https://www.law.cornell.edu/supremecourt/text/22-58}{Cornell LII}; \href{https://www.lawfaremedia.org/article/united-states-v.-texas-explained}{Lawfare}).
\item \textbf{Opposite-Side (Standing Exists):} A state petitions the Environmental Protection Agency (EPA) to initiate a notice-and-comment rulemaking regulating a specific industrial pollutant under the Clean Air Act, and sues when the EPA denies the statutory rulemaking petition, proving imminent coastal erosion to state-owned land. As footnote 6 of \textit{Texas} explains, challenging the denial of a statutorily authorized rulemaking petition is governed by \textit{Massachusetts v. EPA} rather than the non-enforcement standing bar of \textit{Linda R. S.} (\href{https://www.law.cornell.edu/supremecourt/text/22-58}{Cornell LII}).
\item \textbf{Fence-Sitter (Borderline / Unresolved):} Congress amends the INA to add an explicit citizen-suit provision authorizing border states to sue DHS in federal court whenever DHS fails to take custody of covered criminal noncitizens under \S 1226(c) and explicitly empowering courts to issue affirmative arrest injunctions. While Kavanaugh's second caveat in Part II-C suggests that express congressional elevation of injuries and judicial remedies ``might'' change the Article III standing calculus, footnote 4 expressly warns that such a statute ``could also raise Article II issues'' by infringing the President's constitutional enforcement discretion, leaving the outcome sharply contested (\href{https://www.law.cornell.edu/supremecourt/text/22-58}{Cornell LII}).
\end{enumerate}

\section{Critique}

\textbf{Originalist and Textualist Perspectives:} Originalist and textualist scholars offered mixed assessments reflecting the deep split among the conservative Justices themselves. Supporters of Justice Gorsuch's concurrence praised his strict adherence to Article III's case-or-controversy limits, his skepticism of \textit{Massachusetts v. EPA}'s ``special solicitude'' for state plaintiffs, and his textualist challenge to universal vacatur under 5 U.S.C. \S 706(2) (\href{https://www.lawfaremedia.org/article/united-states-v.-texas-explained}{Lawfare}). Conversely, conservative textualists aligned with Justice Alito---as well as commentary in the \textit{Louisiana Law Review}---criticized the majority for nullifying Congress's unambiguous ``shall'' mandates in \S\S 1226(c) and 1231(a)(2) and giving executive agencies effective ``carte blanche'' to disregard statutory commands, undermining both legislative supremacy and the historical meaning of the Take Care Clause (\href{https://lawreview.law.lsu.edu/archives/3081}{Louisiana Law Review}).

\textbf{Progressive and Doctrinal Perspectives:} Doctrinal and progressive scholars generally welcomed the Court's curb on state-attorney-general litigation challenging federal policy through attenuated fiscal-injury theories, which had proliferated in single-judge federal districts (\href{https://www.lawfaremedia.org/article/united-states-v.-texas-explained}{Lawfare}). At the same time, the \textit{Harvard Law Review} critiqued the majority's selective invocation of ``history and tradition,'' arguing that by focusing exclusively on the history of domestic prosecutorial discretion while ignoring the distinct history of state-federal cooperation and criminalization in immigration enforcement, the Court deepened the ``doctrinal fragmentation'' of modern standing law (\href{https://harvardlawreview.org/print/vol-137/united-states-v-texas/}{Harvard Law Review}). Administrative-law scholars also echoed Justice Barrett's concern that constitutionalizing statutory reviewability rules like \textit{Heckler v. Chaney} into Article III standing could unduly restrict Congress's power to authorize judicial checks on executive non-enforcement in future statutory schemes (\href{https://www.law.cornell.edu/supremecourt/text/22-58}{Cornell LII}).

\textbf{Underlying Competing Values:} The case pits \textit{executive flexibility, resource prioritization, and judicial restraint}---preventing federal courts from becoming supervisors of nationwide law-enforcement budgets---against \textit{legislative supremacy, federalism, and the rule of law}, which demand that the Executive Branch faithfully execute mandatory congressional statutes rather than unilaterally suspending them.

\section{Key Quotations}

\begin{quote}
``The States essentially want the Federal Judiciary to order the Executive Branch to alter its arrest policy so as to make more arrests. But this Court has long held `that a citizen lacks standing to contest the policies of the prosecuting authority when he himself is neither prosecuted nor threatened with prosecution.' \textit{Linda R. S. v. Richard D.}, 410 U. S. 614, 619 (1973). Consistent with that fundamental Article III principle, we conclude that the States lack Article III standing to bring this suit.''
\end{quote}
\noindent --- Justice Kavanaugh, Opinion of the Court (\href{https://www.law.cornell.edu/supremecourt/text/22-58}{Cornell LII}).

\begin{quote}
``It bears emphasis that the question of whether the federal courts have jurisdiction under Article III is distinct from the question of whether the Executive Branch is complying with the relevant statutes---here, \S 1226(c) and \S 1231(a)(2). In other words, the question of reviewability is different from the question of legality. We take no position on whether the Executive Branch here is complying with its legal obligations under \S 1226(c) and \S 1231(a)(2). We hold only that the federal courts are not the proper forum to resolve this dispute.''
\end{quote}
\noindent --- Justice Kavanaugh, Opinion of the Court (\href{https://www.law.cornell.edu/supremecourt/text/22-58}{Cornell LII}).

\begin{quote}
``The States lack standing because federal courts do not have authority to redress their injuries. Section 1252(f)(1) denies the States any coercive relief. A vacatur order under \S 706(2) supplies them no effectual relief. And such an order itself may not even be legally permissible. ... The Constitution affords federal courts considerable power, but it does not establish `government by lawsuit.' ''
\end{quote}
\noindent --- Justice Gorsuch, concurring in the judgment (\href{https://www.law.cornell.edu/supremecourt/text/22-58}{Cornell LII}).

\begin{quote}
``The Court weaves together multiple doctrinal strands to create a rule that is not only novel, but also in tension with other decisions. In my view, this case should be resolved on the familiar ground that it must be `likely,' as opposed to merely `speculative,' that any injury will be `redressed by a favorable decision.' ''
\end{quote}
\noindent --- Justice Barrett, concurring in the judgment (\href{https://www.law.cornell.edu/supremecourt/text/22-58}{Cornell LII}).

\begin{quote}
``The majority's conception of Presidential authority smacks of the powers that English monarchs claimed prior to the `Glorious Revolution' of 1688, namely, the power to suspend the operation of existing statutes, and to grant dispensations from compliance with statutes.''
\end{quote}
\noindent --- Justice Alito, dissenting (\href{https://www.law.cornell.edu/supremecourt/text/22-58}{Cornell LII}; \href{https://www.courtlistener.com/opinion/10049663/united-states-v-texas/}{CourtListener}).

\end{document}
